\documentclass[10pt,a4paper]{amsart}
\usepackage[margin=19mm]{geometry}
\usepackage{amsmath,amssymb,mathtools}
\usepackage[scr=rsfs,cal=euler]{mathalfa}
\usepackage{xfrac}
\usepackage[unicode,hidelinks,bookmarks=false]{hyperref}
\newcommand{\ie}{i.e.\ }
\newcommand{\cf}{cf.\ }
\newcommand{\resp}{respectively\ }
\newcommand{\Subsection}{\subsection}
\providecommand{\nobreakdash}{\nobreak\hbox{-}}
\setlength{\emergencystretch}{2em}
\multlinegap0pt
\usepackage{amssymb}
\usepackage{enumerate}
\newtheorem{lemma}{Lemma}
\newtheorem{theorem}{Theorem}
\def\Z{{\mathbb Z}}
\title{Tiling a triangle into a prime number of congruent triangles}
\author{Michael Beeson}
\date{Source: arXiv:2607.23453v1 (CC BY 4.0)}
\begin{document}
\maketitle

\noindent\emph{Source and licence.} Tiling a triangle into a prime number of congruent triangles. Version arXiv:2607.23453v1 (CC BY 4.0), distributed by arXiv under the Creative Commons Attribution 4.0 International licence, http://creativecommons.org/licenses/by/4.0/, as recorded in the arXiv licence metadata for this exact version and shown on the abstract page https://arxiv.org/abs/2607.23453v1. Full licence notice: \texttt{licenses/arxiv-2607.23453-CC-BY-4.0.txt}.\par\smallskip

\noindent\textit{Editorial scope.} Counted unit: the article's theorem theorem:case3 (source0.tex lines 591--596). The article's complete original proof of it is delivered with it (source0.tex lines 598--656, one \texttt{proof} environment of 59 lines). No mathematics is written, repaired or re-derived: the statement and the proof are the article's own lines, in the article's own notation, and no retained line is substituted or re-flowed.  Retained uncounted as the source-local context the proof needs: lines 380--390; lines 413--423; lines 432--434; lines 449--454; lines 494--496; lines 513--517.  Nothing was omitted from any retained span, so the package carries no omission record.  No retained line contains a citation, so the wrapper carries no bibliography and no bibliography entry.  No retained line contains a figure environment, an image-inclusion command or drawing code, so no image is delivered and the package declares no asset dependency.  Editorial formatting only, in addition: an AMS \texttt{amsart} preamble that restates the article's own declarations and macros used by the retained lines, namely the environments \texttt{lemma}, \texttt{theorem} on separate counters, together with the standard packages those lines need.  The retained environments print in this document as Lemma 1, Theorem 1 in order of appearance, whereas the article prints the same items with the numbering of its own full document.  The complete native source of the article as distributed in the arXiv e-print for this exact version is bundled unchanged as \texttt{sources/206029/source0.tex}.
\begin{lemma} \label{lemma:claude}
Let triangle $R$ have angles $(\alpha,\beta,2\pi/3)$ and sides $(a,b,c)$ opposite
those angles, respectively. Then
\begin{eqnarray*}
\cos \alpha &=& \frac{a+2b} {2c} \\
\sin \alpha &=& \frac {a \sqrt 3}{2c}\\
\sin 2 \alpha &=& \frac {a \sqrt 3 (a+2b)}{2c^2}\\
\sin 3 \alpha 
&=&  \frac {a \sqrt 3}{2c} \cdot\frac {3b(a+b)}{c^2}
\end{eqnarray*}
\end{lemma}

\begin{lemma}\label{lemma:beta}
Let triangle $R$ have angles $(\alpha,\beta,2\pi/3)$ and sides $(a,b,c)$ opposite
those angles, respectively. Then
\begin{eqnarray*}
\sin \beta &=& \frac{b \sqrt 3}{2c}\\
\cos \beta &=& \frac {2a+b}{2c}\\
\sin 2 \beta &=& \frac {b \sqrt 3 (2a+b)}{2c^2}\\
\sin (2 \alpha +\beta) &=& \frac{\sqrt 3 (a+b)}{2c} \\
\sin (\alpha + 2 \beta)&=&\frac{\sqrt 3 (a+b)}{2c} 
\end{eqnarray*}
\end{lemma}

\begin{lemma} \label{lemma:pairwisecoprimality} Suppose $c^2 = a^2 + b^2 + ab$ and 
$(a,b,c)$ have no common divisor.  Then $(a,b,c)$ are pairwise relatively prime.
\end{lemma}

\begin{lemma} \label{lemma:claude3}
Suppose triangle $T$ is tiled by a triangle $R$ with integer sides.
 Let the side lengths 
$(X,Y,Z)$ of $T$ be proportional to coprime integers $(X_0, Y_0, Z_0)$.
Then $(X,Y,Z) = \lambda (X_0,Y_0,Z_0)$ for some integer $\lambda$.
\end{lemma}

\begin{lemma} \label{lemma:not3divc} Let $(a,b,c)$ have no common factor
and satisfy $c^2 = a^2 + ab + b^2$.  Then $3 \nmid c$.
\end{lemma}

\begin{lemma}\label{lemma:rationality3}
Suppose triangle $T$ with incommensurable angles is tiled using triangle $R$ with angles $(\alpha, \beta, 2\pi/3)$.
Suppose $T$ is not similar to $R$.
Then the side lengths $(a,b,c)$ of $R$ are commensurable.
\end{lemma}

\begin{theorem} \label{theorem:case3}  Let triangle $T$ have incommensurable angles 
$(\alpha, 2\beta, 2\alpha + \beta)$, and suppose $T$ is $N$-tiled by $(\alpha,\beta,2\pi/3)$.
Then $N$ is not a prime.  Moreover,
$$ N = \lambda^2 (2a+b) (a+b),$$
where $(a,b,c)$ are the sides of the tile with no common factor, and $\lambda \in \Z$.
\end{theorem}

\begin{proof}  Since $\alpha, \beta < \pi/3$,  neither $\alpha$ nor $2\beta$ is equal to $2\pi/3$.
If $2\alpha + \beta = 2\pi/3$ then $\alpha = \pi/3$, since $\alpha + \beta = \pi/3$; but that
violates the incommensurable angles hypothesis.  Hence $T$ has no $2\pi/3$ angle. 
Therefore $T$ is not similar to the tile.
By Lemma~\ref{lemma:rationality3}, $(a,b,c)$ are commensurable.
Without loss of generality, we may assume they are integers with no common factor.
Then the side lengths $(X,Y,Z)$ are also integers, since they are integral linear
combinations of $(a,b,c)$. 
Let the sides of $T$ next to the $2\alpha + \beta$ angle be $X$ and $Y$.  
We will show that $(X,Y,Z)$  are proportional to a triple with
no common factor, apply Lemma~\ref{lemma:claude3} to conclude
that the proportionality factor is an integer, 
and then use the area equation to finish.   Here are the details:

By the law of sines,
$$ X: {\sin \alpha} =   Y :{\sin 2 \beta} = Z :{\sin( 2\alpha + \beta)}.$$
Using the formulas from Lemmas~\ref{lemma:claude} and~\ref{lemma:beta}, we have 

$$ X: \frac {a \sqrt 3}{2c} = Y:\frac {b \sqrt 3 (2a+b)}{2c^2} = Z: \frac{\sqrt 3 (a+b)}{2c}.$$
$$ X: ac = Y : b(2a+b) = Z: (a+b)c.$$
Then for some real $\lambda$ we have
\begin{equation}
 (X,Y,Z) = \lambda (ac, b(2a+b), (a+b)c). \label{eq:518}
 \end{equation}
We will prove $(ac,b(2a+b),(a+b)c)$ have no common factor. 
We have $c^2 = a^2 + b^2 + ab$; by Lemma~\ref{lemma:pairwisecoprimality},
we have $\gcd(a,b) = 1$. 
Suppose the prime $p$ divides all three of $(ac,b(2a+b),(a+b)c)$.
  Then $p | ac$, so 
$p | a$ or $p | c$.  Then $p \nmid b$.  But $p | b(2a+b)$,
so $p | (2a+b)$.  If $p | a$ then $p | b$; contradicting $\gcd(a,b) = 1$. 
Therefore $p \nmid a$. Therefore $p | c$, since $p | a$ or $p | c$. 
Modulo $p$ we have
\begin{eqnarray*}
b \equiv -2a   && \mbox{\qquad since $p | (2a+b)$}\\
0 \equiv c^2 \equiv a^2 + b^2 + ab && \mbox{\qquad since $p | c$}\\
0  \equiv 3a^2 \pmod p  && \mbox{\qquad since $b \equiv -2a$}\\
p = 3                && \mbox{\qquad since $p \nmid a$}\\
\end{eqnarray*}
But that contradicts Lemma~\ref{lemma:not3divc}.
Therefore $(ac,b(2a+b),(a+b)c)$ have no common factor. 
Then by Lemma~\ref{lemma:claude3}, $\lambda$ is a positive integer in (\ref{eq:518}).

The area equation for this shape of triangle is 
$$ XY \sin (2\alpha+\beta) \ = \ N ab \frac {\sqrt 3} 2.$$
Substituting the values for $X$ and $Y$ from (\ref{eq:518}) into the area equation,
we obtain
\begin{eqnarray*}
\lambda^2 abc(2a+b) \sin(2\alpha+\beta) &=& \ N ab \frac {\sqrt 3} 2  \\
\lambda^2 c(2a+b) \sin(2\alpha+\beta)  &=& N \frac {\sqrt 3} 2  
\end{eqnarray*}
Putting in the value of $\sin(2\alpha+\beta)$ from Lemma~\ref{lemma:beta}, we have
\begin{eqnarray*}
\lambda^2 c(2a+b) \sqrt 3 \frac {a+b}{2c}  &=& N \frac {\sqrt 3} 2  \\
\lambda^2 (2a+b) (a+b)  &=& N 
\end{eqnarray*}
Since $\lambda$ is an integer, and $a,b,c$ are positive, we have a factorization of $N$.
Hence $N$ is not prime.
\end{proof}

\end{document}
